\section{Related Work}
\paragraph{Hypergraph Construction}
 Existing hypergraph construction methods are mostly based on neighbor reconstruction and neighbor distances. For example, the $k$ nearest neighbor-based method \citep{huang2009video} connects a centroid node and its $k$ nearest neighbors in the feature space to form a hyperedge. \citet{wang2015visual,liu2017elastic,jin2019robust,xiao2019multi} further refine these neighbor connections through various regularization. However, the number of hyperedges of these methods is restricted to the number of nodes, and hyperedges obtained by these methods are inconsistent across instances. 
 Cluster-based approaches learn community structure in a graph \cite{bannadabhavi2023community,ying2018hierarchical}, which can also be used to form hyperedges.
 \citet{zhang2022deep,zhang2018dynamic} proposed to iteratively refine a noisy hypergraph, which is obtained by the aforementioned methods. Therefore, they share the same limitations as the aforementioned methods. In addition, these methods are unable to learn MIMR hyperedges due to the absence of a tractable objective. Other methods, such as attributed-based methods~\citep{huang2015learning, joslyn2019high}, are ill-adapted to our context since they require discrete labels. 
 %or a prior graph topology, which is usually very noisy.
 Different from these methods, we provide a way to learn a consistent hypergraph through a deep-learning model. Furthermore, thanks to the proposed \multiheadname, these hyperedges are theoretically ensured MIMR.

\paragraph{High-Order Relationships in fMRI}

Although there are some methods working on high-order relationships in fMRI, they have limitations and are inconsistent with our MIMR objective.
\citet{xiao2019multi,li2022construction} used the existing non-learning-based hypergraph construction methods, which may lead to noisy and inexpressive hypergraphs. \citet{zu2016identifying,santoro2023higher} enumerated all hyperedges with degrees lower than 3, which can only discover a tiny portion of all possible hyperedges in exponential space and is not scalable to a large degree. \citet{rosas2019quantifying} proposed O-information, which reflects the balance between redundancy and synergy. The O-information metric is utilized by \citet{varley2023multivariate} to study fMRI data. However, the objective of these methods is not consistent with ours: although both of us are quantifying the redundancy of high-order relations, our method is to learn those that are most informative toward a cognition score, while theirs is to depict the synergy and redundancy within a system.

\paragraph{Information Bottleneck}
Information bottleneck (IB) \citep{tishby2000information} is a technique in data compression. The key idea is to extract a summary of data, which contains the most relevant information to the objective. \citet{alemi2016deep} first employed an IB view of deep learning. After that, IB has been widely used in deep learning. The applications span areas such as computer vision \citep{luo2019significance,peng2018variational}, reinforcement learning \citep{goyal2019infobot,igl2019generalization}, natural language processing \citep{wang2020learning} and graph learning \citep{yu2020graph,yu2022improving, xu2021infogcl,wu2020graph}. Unlike these studies that use IB to extract a compressed representation or a select set of features, our approach focuses on identifying the underlying structures of the data.

\paragraph{Connectivity-based Phenotypic Prediction}
Recently, deep learning techniques have been increasingly employed in predicting phenotypic outcomes based on the connectivity of brain regions. Most works \citep{ahmedt2021graph,li2019graph,cui2022interpretable,kan2022fbnetgen,cui2022braingb,said2023neurograph} model the brain network as a graph, in which regions act as nodes and pairwise correlations form the edges. These methods predominantly utilize Graph Neural Networks (GNNs) to capture the connectivity information for predictions. In addition to GNNs, \citet{kan2022brain} proposed to use transformers with a specially designed readout module, leveraging multi-head attention mechanisms to capture pairwise connectivity. However, all of these methods heavily rely on pairwise connectivity and neglect more intricate higher-order relationships. This oversight, on the one hand, leads to sub-optimal prediction performances and, on the other hand, prevents domain experts from acquiring insightful neuroscience interpretations, given that brain functions often involves multiple regions.